\section{Study Design}\label{sec:design}

\subsection{System under test}
The service answers a question in three steps. (1) \emph{Retrieval}: dense retrieval with
\texttt{bge-small-en-v1.5}~\cite{xiao2023cpack} over a fixed passage store, top $k{=}5$. A BM25 index
over the same store is available as a lexical fallback. (2) \emph{Prompt assembly}: each passage is
rendered as \texttt{[P$i$] title (last updated date)} followed by its text, and the instruction asks
for the shortest possible answer plus the ids of supporting passages, or \INS{}.
(3) \emph{Serving}: a 4-bit instruction-tuned model served locally with MLX~\cite{mlx2023}, greedy
decoding. The service logs a trace for every call: retrieved documents with their similarity scores
and metadata, the passages exactly as rendered, the number of prompt tokens it sent and the number
the server reports as processed, and the raw output. Repair policies see this trace and nothing else;
they never see the fault label or the gold answer.

\subsection{Data}
We use the HotpotQA distractor validation set~\cite{yang2018hotpotqa}. We keep \emph{bridge}
questions whose answer string occurs in exactly one of the question's ten paragraphs, in a gold
supporting sentence, and drop answers with annotation artifacts. The first 1{,}200 such questions
(seeded shuffle) define the store: all their paragraphs, deduplicated by title (11{,}737 passages).
Of these, 664 have both gold paragraphs in the clean dense top-5; this filter guarantees that in the
clean condition the evidence is present, so each injected fault is the only reason it goes missing.
We draw 30 development and 150 test questions and exclude every question used in an earlier
prototype of this study. Each passage carries a ``last updated'' date (2023--2024). For the stale
fault, each question's answer passage also has an \emph{old version}, dated three to six years
earlier, in which the answer value is replaced by a plausible wrong value: numbers and dates are
shifted by rule; entity answers are rewritten by a local 7B model (e.g., Sam Waterston $\to$ Tom
Hanks) and checked to differ from the gold answer.

\subsection{Faults}
Table~\ref{tab:faults} lists the six faults. Each models a failure reported for deployed RAG
services, and none targets the answer passage directly except \emph{absent}, where abstaining is the
correct behavior.

\begin{table}[t]
\caption{Injected faults, by pipeline stage.}\label{tab:faults}
\small
\begin{tabular}{@{}p{1.55cm}p{3.15cm}p{2.9cm}@{}}
\toprule
Fault & Injection & Production cause \\
\midrule
encoder & queries embedded with \texttt{all-MiniLM-L6-v2}~\cite{reimers2019sbert}; index built with \texttt{bge-small} & query service falls back to a library default encoder \\
truncation & template cuts every passage to its first 100 characters & \texttt{text[:max\_chars]} misconfiguration \\
overflow & server keeps only the last 512 prompt tokens & unset context window, e.g.\ \texttt{num\_ctx} default~\cite{musetoiu2026ragal,ragflow_ollama_ctx} \\
stale & old version of the answer passage indexed beside the current one; retrieval ranks both & index not deduplicated after an update~\cite{ouyang2025hoh} \\
absent & answer passage removed from the index & content gap \\
enc.+trunc. & encoder and truncation together & independent faults co-occur \\
\bottomrule
\end{tabular}
\end{table}

\subsection{Signals and repair policies}
\emph{Trace checks} are computed from the trace with no model call: (i) \emph{truncation}, a rendered
passage is shorter than the retrieved text; (ii) \emph{version conflict}, two passages share a title;
(iii) \emph{overflow}, the server processed fewer prompt tokens than were sent; (iv) \emph{low
retrieval}, the top dense similarity is below the 5th percentile of clean traffic, calibrated on
467 held-out questions. Each check has a matching repair: render in full, keep the newest version,
fit the passages to the window the server reported, or fall back to BM25. \emph{Support checks}
decide whether an answer is grounded: a text-match check (the answer appears in the passages), the
model's own verifier (``Is the proposed answer correct and directly supported by the passages?
SUPPORTED/UNSUPPORTED''), or MiniCheck-RoBERTa-Large~\cite{tang2024minicheck}, a 355M-parameter
grounding classifier run on the passages packed into 400-token chunks. MiniCheck's model and claim
format were chosen on prototype development data; its thresholds were set on the development split by
rules fixed in advance (0.5, MiniCheck's default, to trigger a repair; the 2nd percentile of scores
of correct grounded answers, 0.01, to abstain). The support check of the proposed policy (text
match) was also chosen on the development split; the other two remain as comparison arms.

Table~\ref{tab:policies} lists the policies. All start from the same initial answer and may make at
most two more model calls; LLM verifier calls count against this budget, MiniCheck calls do not
(their time is reported separately). \emph{Diagnose} applies the repairs indicated by the trace
checks, then checks support; while the answer is unsupported and budget remains it widens
retrieval to 10, then adds BM25 results, regenerating after each step; if the final answer is
still clearly unsupported it returns \INS{}. The self-check and diagnose variants
form a factorial design: \{LLM verifier, MiniCheck\} $\times$ \{regenerate on the same passages,
repair the context\}, plus text match with context repair.

\begin{table}[t]
\caption{Repair policies (budget: two extra model calls).}\label{tab:policies}
\small
\begin{tabular}{@{}p{2.25cm}p{5.4cm}@{}}
\toprule
Policy & Behavior \\
\midrule
No repair & deployed pipeline \\
Retry & if \INS{}, resample at temperature 0.7 \\
Self-check (2 verifiers) & verify the answer (LLM or MiniCheck); if unsupported, regenerate on the same passages with a hint \\
Repair-all & always: dense top-10 $\cup$ BM25 top-5, full rendering, newest versions, regenerate \\
Trace checks only & trace checks and their repairs; no support check \\
Diagnose (3 support checks) & trace checks, then a support check (text match, LLM, or MiniCheck) with re-check after each repair, retrieval ladder, abstention \\
\quad ablations & proposed (text match) variant without re-check, without abstention, without trace checks \\
\bottomrule
\end{tabular}
\end{table}

\subsection{Measures and statistics}
\emph{Success} is a correct answer in the answerable conditions and \INS{} in
\emph{absent}. Correctness is exact match after normalization (number words mapped to digits),
token F1 $\geq 0.8$, or the gold phrase inside a short prediction; other predictions go to a local
answer-equivalence judge (Qwen2.5-7B) that sees only the question, gold answer and prediction. On
the prototype, 197 of 203 audited judge decisions agreed with a manual audit; we report results with
string matching alone as a sensitivity check. \emph{Success under faults} averages the five
answerable fault conditions per question. \emph{Clean damage} is the share of clean questions
answered correctly at first and wrongly after the policy~\cite{dutta2024accuracy}. Costs are model
calls, prompt tokens and GPU seconds per case. All comparisons are paired by question: 95\% intervals
are percentile bootstraps over questions (2{,}000 resamples), per-condition tests are exact McNemar.

\subsection{Model and setup}
The generator is Llama-3.2-3B-Instruct~\cite{grattafiori2024llama3}, 4-bit, served with MLX on one
MacBook Pro (M4, 16\,GB); this paper reports one model, and extending the study to
Qwen2.5-7B~\cite{qwen2025qwen25} and Gemma-2-9B~\cite{gemma2024gemma2} is ongoing
(Section~\ref{sec:threats}). A response cache keyed by model, prompt and sampling settings lets
policies that issue the same prompt share generations and makes every run replayable.
